% Configuration PDF commune aux documents Husson.
\usepackage{fontspec}
\IfFontExistsTF{texgyretermes-regular.otf}{%
  \defaultfontfeatures{%
    SmallCapsFont = texgyretermes-regular.otf,%
    SmallCapsFeatures = {Letters=SmallCaps}%
  }%
}{%
  \defaultfontfeatures{SmallCapsFeatures={Letters=SmallCaps}}%
}
\usepackage{xcolor}
\usepackage{fvextra}
\usepackage{newunicodechar}
\usepackage{enumitem}
\usepackage{etoolbox}
\usepackage{needspace}
\usepackage{pdflscape}
\usepackage{adjustbox}

\definecolor{headingblue}{HTML}{17365D}
\definecolor{headingred}{HTML}{7A1F2B}
\definecolor{headinggreen}{HTML}{1F5A3A}
\definecolor{headingcharcoal}{HTML}{34383D}
\definecolor{commentgray}{HTML}{6F6F6F}
\definecolor{pythonaccent}{HTML}{9A6100}
\definecolor{pythonfill}{HTML}{FFF5D9}
\definecolor{raccent}{HTML}{276DC3}
\definecolor{rfill}{HTML}{EAF2FB}

\newunicodechar{≥}{\ensuremath{\geq}}
\newunicodechar{≤}{\ensuremath{\leq}}
\setlength{\emergencystretch}{3em}

% Pagination éditoriale. Une ligne isolée au début ou à la fin d'une page est
% évitée, tandis que \raggedbottom autorise un blanc final raisonnable plutôt
% que d'étirer artificiellement les espacements verticaux.
\clubpenalty=10000
\widowpenalty=10000
\displaywidowpenalty=10000
\raggedbottom

\makeatletter
\renewcommand*{\@pnumwidth}{2.8em}
\renewcommand*{\@tocrmarg}{3.6em}
\makeatother

% Hiérarchie éditoriale sobre : les trois premiers niveaux structurent la page
% par la graisse et une couleur foncée. Les niveaux suivants restent proches
% du texte courant et se distinguent par la graisse, l'italique ou les petites
% capitales. Seuls les trois premiers niveaux sont numérotés.
\setcounter{secnumdepth}{3}
\setkomafont{disposition}{\rmfamily}
\setkomafont{section}{\rmfamily\color{headingblue}\bfseries\Large}
\setkomafont{subsection}{\rmfamily\color{headingred}\bfseries\large}
\setkomafont{subsubsection}{\rmfamily\color{headinggreen}\bfseries\normalsize}
\setkomafont{paragraph}{\rmfamily\color{headingcharcoal}\bfseries\normalsize}
\setkomafont{subparagraph}{\rmfamily\color{headingcharcoal}\itshape\normalsize}

\RedeclareSectionCommand[
  beforeskip=1.25\baselineskip,
  afterskip=0.65\baselineskip
]{section}
\RedeclareSectionCommand[
  beforeskip=1.05\baselineskip,
  afterskip=0.50\baselineskip
]{subsection}
\RedeclareSectionCommand[
  beforeskip=0.90\baselineskip,
  afterskip=0.42\baselineskip
]{subsubsection}
\RedeclareSectionCommand[
  beforeskip=0.76\baselineskip,
  afterskip=0.28\baselineskip
]{paragraph}
\RedeclareSectionCommand[
  beforeskip=0.62\baselineskip,
  afterskip=0.24\baselineskip
]{subparagraph}

% Pandoc rend normalement un titre Markdown de niveau 6 comme du texte brut
% en LaTeX. Le filtre heading-levels.lua l'entoure de cet environnement pour
% lui donner un style discret et un espace propre avec le texte.
\newenvironment{hussonsixthheading}{%
  \Needspace{3\baselineskip}%
  \par\addvspace{0.52\baselineskip}%
  \noindent\rmfamily\color{headingcharcoal}\scshape\small
}{%
  \par\nobreak\vspace{0.20\baselineskip}%
}

\setlength{\parindent}{0pt}
\setlength{\parskip}{3pt}
\setlist[itemize]{itemsep=1pt,parsep=1pt,topsep=2pt}
\setlist[enumerate]{itemsep=1pt,parsep=1pt,topsep=2pt}

% Espacement vertical compact et sobre pour les tableaux (hauteur minimale).
% La boîte typographique normale est rétablie avant de composer les paragraphes
% multiligne : la marge n'est donc appliquée qu'une fois par cellule.
\renewcommand{\arraystretch}{1}
\setlength{\extrarowheight}{0pt}
\makeatletter
\newbox\hussontablestrutbox
\newcommand{\hussontablepadding}{%
  \setbox\hussontablestrutbox=\copy\strutbox
  \setbox\strutbox=\hbox{%
    \vrule width 0pt
      height \dimexpr\ht\hussontablestrutbox+2.5pt\relax
      depth \dimexpr\dp\hussontablestrutbox+2.5pt\relax
  }%
  \let\hussontablesavedmkpream\@mkpream
  \def\@mkpream##1{%
    \hussontablesavedmkpream{##1}%
    \setbox\strutbox=\copy\hussontablestrutbox
    \let\@mkpream\hussontablesavedmkpream
  }%
}
\makeatother

% Les tableaux gardent leur hauteur naturelle, avec une taille homogène. Un
% tableau long peut se poursuivre sur la page suivante, mais ne commence pas
% dans les toutes dernières lignes de la page courante.
\AtBeginEnvironment{longtable}{\Needspace{6\baselineskip}\small\hussontablepadding}
\AtBeginEnvironment{tabular}{\small\hussontablepadding}

% Filet de securite : seuls les tableaux depassant la ligne sont reduits a \linewidth.
\BeforeBeginEnvironment{tabular*}{\begin{adjustbox}{max width=\linewidth}}
\AfterEndEnvironment{tabular*}{\end{adjustbox}}
\BeforeBeginEnvironment{tabular}{\begin{adjustbox}{max width=\linewidth}}
\AfterEndEnvironment{tabular}{\end{adjustbox}}

% Les blocs de code restent divisibles lorsqu'ils sont longs. Cette réserve
% empêche seulement qu'un bloc commence en bas de page avec une ou deux lignes.
\BeforeBeginEnvironment{Shaded}{\Needspace{8\baselineskip}}

% Les blocs de code sont compacts et peuvent revenir à la ligne.
\DefineVerbatimEnvironment{Highlighting}{Verbatim}{%
  fontsize=\footnotesize,
  breaklines,
  breakanywhere,
  commandchars=\\\{\}
}
\providecommand{\CommentTok}[1]{\textcolor{commentgray}{\textit{#1}}}

\newcommand{\codelanguagelabel}[2]{%
  \Needspace{5\baselineskip}%
  \par\addvspace{0.55\baselineskip}%
  \noindent\begingroup
  \setlength{\fboxsep}{3pt}%
  \colorbox{#1}{\color{white}\sffamily\bfseries\scriptsize #2}%
  \endgroup\par\nobreak\vspace{-0.25\baselineskip}%
}

\newenvironment{pythoncode}{%
  \codelanguagelabel{pythonaccent}{PYTHON}%
  \begingroup\colorlet{shadecolor}{pythonfill}%
}{%
  \endgroup\par\addvspace{0.25\baselineskip}%
}

\newenvironment{rcode}{%
  \codelanguagelabel{raccent}{R}%
  \begingroup\colorlet{shadecolor}{rfill}%
}{%
  \endgroup\par\addvspace{0.25\baselineskip}%
}

\makeatletter
\renewcommand{\maketitle}{
  \begin{center}
    {\Large\bfseries\rmfamily \@title\par}
    \vskip 0.35em
    {\large\rmfamily \@subtitle\par}
    \vskip 0.65em
    {\normalsize\rmfamily \@author\par}
  \end{center}
}
\makeatother

\renewcommand{\contentsname}{}
\AtBeginDocument{
  \let\oldtableofcontents\tableofcontents
  \renewcommand{\tableofcontents}{
    \begingroup
    \footnotesize
    \setlength{\parskip}{1pt}
    \oldtableofcontents
    \endgroup
  }
}
